
\section{\trBrokenTrajectory{}与\trGluing：陈述}
\label{sec:broken-gluing}

在\cref{subsec:broken-trajectories}中，我们已经证明：把两个\trCriticalPoint{} $a$ 与 $b$ 连接起来的\trFloerEquation{}诸解的一个序列可能\trConvergence{}到一条“\trBrokenTrajectory”。破碎\trTrajectory{}所成的空间表现为\trTrajectory{}空间 $\scr{L}(a,b)$ 的一个“\trCompactification”。让我们在 $a$ 与 $b$ 的\trIndex{}之差为 $2$ 的情形下，更仔细地研究它的结构。

\begin{theorem}\label{thm:lbar-manifold-boundary}
设 $(H,J)$ 是一个\trRegular{}偶对（其中 $H$ 是\trNondegenerate{}的）。设 $x$ 与 $z$ 是 $H$ 的两条\trPeriodic\trTrajectory，它们的\trIndex{}满足
\[
\mu(x)=\mu(z)+2\text{。}
\]
那么 $\overline{\scr{L}}(x,z)$ 是一个维数为 $1$ 的\trCompact\trManifoldWithBoundary，并且\footnote{译注：原书此式左端印作 $\partial\overline{\scr{L}}(x,y)$；按本定理主句所断言的带边\trManifold{}为 $\overline{\scr{L}}(x,z)$，疑应为 $\partial\overline{\scr{L}}(x,z)$。又，右端并集的下标印作 $\mu(x)<\mu(y)<\mu(z)$；按假设 $\mu(x)=\mu(z)+2$ 以及诸\trTrajectory{}沿\trIndex{}递降的口径，疑应为 $\mu(z)<\mu(y)<\mu(x)$。此处均径改。——译者注}
\[
\partial\overline{\scr{L}}(x,z)=\bigcup_{\mu(z)<\mu(y)<\mu(x)}\scr{L}(x,y)\times\scr{L}(y,z)\text{。}
\]
\end{theorem}

\begin{corollary}\label{thm:boundary-square-zero}
在同样的假设下，
\[
\partial\circ\partial=0\text{。}
\]
\noqed
\end{corollary}

本章的其余部分都用于\cref{thm:lbar-manifold-boundary}的证明。我们已经在\cref{subsec:broken-trajectories}中证明了：\trTrajectory{}空间 $\overline{\scr{L}}(x,z)$ 是\trCompact{}的，而 $\scr{L}(x,z)$ 是一个维数为 $1$ 的\trManifold。因此，只需研究在边界点的邻域中所发生的事情。这将在下述所谓的“\trGluing”定理中得到描述。

\begin{theorem}\label{thm:gluing-psi}
设 $x$、$y$、$z$ 是\trActionFunctional{} $\scr{A}_{H}$ 的三个具有相继\trIndex{}的\trCriticalPoint，
\[
\mu(x)=\mu(y)+1=\mu(z)+2\text{，}
\]
并设 $(u,v)\in\scr{M}(x,y)\times\scr{M}(y,z)$ 代表诸\trTrajectory
\[
(\widetilde{u},\widetilde{v})\in\scr{L}(x,y)\times\scr{L}(y,z)\text{。}
\]
那么我们有：
\begin{itemize}
\item 存在一个可微映射 $\psi\colon[\rho_{0},+\infty[\to\scr{M}(x,z)$（对某个 $\rho_{0}>0$），使得 $\pi\circ\psi$ 是一个\trEmbedding
\[
\widehat{\psi}=\pi\circ\psi\colon[\rho_{0},+\infty[\longrightarrow\scr{L}(x,z)
\]
并且满足
\[
\lim_{\rho\to+\infty}\widehat{\psi}(\rho)=(\widetilde{u},\widetilde{v})\in\overline{\scr{L}}(x,z)\supset\scr{L}(x,z)\text{。}
\]
\item 此外，若 $\ell_{n}\in\scr{L}(x,z)$ 是一个趋于 $(\widetilde{u},\widetilde{v})$ 的序列，那么对充分大的 $n$，$\ell_{n}\in\mathrm{Im}(\widehat{\psi})$。
\end{itemize}
\end{theorem}

\begin{remark}\label{thm:last-assertion-essential}
最后这条断言与它在\trMorse{}情形中的类似断言（见\cref{thm:single-branch}）同样本质。
\end{remark}

这个证明比起\trMorse{}理论中对应结果（即\cref{thm:gluing-embedding}）的证明来，要困难得多。我们将把它分解成三个步骤：

（1）一个“\textbf{\trPregluing}”（pre-gluing）：我们在其中构造一个\textbf{\trInterpolation}（interpolation） $w_{\rho}$，它依赖于一个参数 $\rho$，介于\trFloerEquation{}的两个解 $u$ 与 $v$ 之间；它将不完全是一个解，而将是（按我们稍后将要给出的意义而言的）一个“近似”解。

（2）$\psi$ 的构造，我们将把它写成
\[
\psi(\rho)=\exp_{w_{\rho}}\bigl(\gamma(\rho)\bigr)
\]
其中
\[
\gamma(\rho)\in W^{1,p}\bigl(w_{\rho}^{\star}TW\bigr)=T_{w_{\rho}}\scr{P}(x,z)
\]
使得 $\psi(\rho)$ 是\trFloerEquation{}的一个真正的解；我们将从 $w_{\rho}$ 出发、经由一个归功于\trPicard{}的\trNewton{}法变体来得到它，我们把这个变体称作“\trNewton--\trPicard{}方法”。

（3）验证 $\widehat{\psi}=\pi\circ\psi$ 的三条所需要的性质，即：
\begin{itemize}
\item $\lim_{\rho\to+\infty}\widehat{\psi}(\rho)=(\widetilde{u},\widetilde{v})$；
\item $\widehat{\psi}$ 是一个\trEmbedding；
\item $\widehat{\psi}$ 在\cref{thm:gluing-psi}最后一条断言的意义下是唯一的。
\end{itemize}

为了这些证明，我们将需要关于\trExponentialMap{}的各阶\trDifferential{}以及\trFloer{}算子的种种估计。它们的证明并不困难，但是却非常冗长（而且有些读者可能会发现它们“精神上显然如此”）。出于这个原因，也为了减轻阅读的负担，我们将把它们集中在第 13 章里给出。当这些技术性引理在正文中出现时，我们会随时指出。
